% Programming with InitechBase, Chapter 19: Memo Fields in Programs (AN-16).
% Panels replayed (examples/p19-*.dot); sources fs/dbt.c (the memo text
% read with its record), core/fn_builtins.c (memo accepted where <expC> is),
% the program NOTES.PRG.  A memo field is read by programs; REPLACE into a
% memo field does not store in this version (Using InitechBase, REPLACE).
\chapter{Memo Fields in Programs}
\Gates{test-pubs-examples, test-dbt-read, test-xbase-fn-a}
The field \texttt{NOTES} of the staff roster is a memo field: its text, of any
length, is kept in \file{STAFF.DBT}.  A program reads it like a character
field.

\begin{ChapterCovers}
\item Memo text as a character value
\item Functions on memo fields
\item Printing the remarks of the roster
\end{ChapterCovers}

\ATEHead{Preparing for This Chapter}
Start InitechBase with a fresh copy of the sample diskette, and type
\type{USE STAFF}.

\ATEHead{Memo Text}
In an expression, a memo field is the text of the current record's memo.
\fn{TYPE()} reports it as M; the string functions accept it wherever they
accept a character value:
\DotExample{examples/p19-memo}
A record with no memo has an empty memo field.

\ATEHead{The Remarks of the Roster}
\file{NOTES.PRG} prints the remark of every employee who has one:
\ProgramListing{NOTES.PRG}
\DotExample{examples/p19-notes}
\fn{TRIM()} and \fn{LEN()} together tell whether the memo holds any text.

\begin{AteNote}
The programs of this manual read memo fields; they do not change them.
\end{AteNote}

\ATEHead{Where Do You Go From Here?}
Chapter~20 changes the order of the records and hides some of them.

\ATEHead{Summary}
\begin{enumerate}
\item A memo field is used as a character value in expressions.
\item \fn{TYPE()} of a memo field is M.
\item \fn{LEN(TRIM(NOTES))} is 0 for a record without a remark.
\end{enumerate}
